% Generated draft: requires mathematical and source review.
\documentclass[11pt]{article}
\usepackage[margin=25mm]{geometry}
\usepackage{fontspec}
\setmainfont{Latin Modern Roman}
\usepackage{amsmath,amssymb,mathtools,unicode-math}
\setmathfont{Latin Modern Math}
\usepackage{enumitem,needspace,xurl,hyperref}
\setlength{\parindent}{0pt}
\setlength{\parskip}{5pt}
\setlength{\emergencystretch}{3em}
\newcommand{\N}{\mathbb{N}}
\newcommand{\Z}{\mathbb{Z}}
\newcommand{\Q}{\mathbb{Q}}
\newcommand{\R}{\mathbb{R}}
\newcommand{\C}{\mathbb{C}}
\newcommand{\F}{\mathbb{F}}
\newcommand{\A}{\mathbb{A}}
\newcommand{\PP}{\mathbb{P}}
\newcommand{\E}{\mathbb{E}}
\newcommand{\eps}{\varepsilon}
\newcommand{\dd}{\,\mathrm{d}}
\begin{document}
\section*{UnsolvedMath ID: 30002885}
\textbf{Source problem code:} catalogue entry
\subsection{Definitions and mathematical statement}
Let \(G\) be a semisimple complex group, let \(G^\vee\) be its Langlands dual, and write \(\mathfrak g=\operatorname{Lie}(G)\). For a complex algebraic group \(H\), the usual formal-loop description of its affine Grassmannian is
\[
\operatorname{Gr}_H=H(\C((t)))/H(\C[[t]]),
\]
where \(\C[[t]]\) and \(\C((t))\) denote the formal power series and Laurent series rings. This description explains the notation but does not resolve whether the source intends \(\mathrm{Gr}=\operatorname{Gr}_G\) or \(\mathrm{Gr}=\operatorname{Gr}_{G^\vee}\): the record introduces both groups and says only ``its affine Grassmannian.''

Write \(K\mathrm{Conv}(\mathrm{Gr})\) for the source's K-theoretic convolution category and \(\mathcal M_q(G)\) for its intended category of \(U_q(\mathfrak g)\)-equivariant \(\mathcal O_q(G)\)-modules. Here \(U_q(\mathfrak g)\) and \(\mathcal O_q(G)\) denote the intended quantum enveloping algebra and quantum coordinate algebra. Their precise versions, the coefficient ring or field, the allowed values or formal meaning of \(q\), and the compatibility defining equivariance are not specified in the record.

\textbf{Conjectural statement supplied by the source.} For every semisimple complex group \(G\), with these constructions interpreted as intended in the source, there should exist a full subcategory \(\mathcal A\subseteq\mathcal M_q(G)\) and an equivalence of categories
\[
K\mathrm{Conv}(\mathrm{Gr})\simeq\mathcal A.
\]
Fullness means that, for all objects \(M,N\) of \(\mathcal A\),
\[
\operatorname{Hom}_{\mathcal A}(M,N)
=\operatorname{Hom}_{\mathcal M_q(G)}(M,N).
\]

\textbf{Unresolved qualifications.} The record does not define the objects, morphisms, equivariance, coefficient choices, or convolution composition of \(K\mathrm{Conv}(\mathrm{Gr})\), nor does it characterize the objects of the proposed full subcategory. It also does not specify whether the equivalence must preserve additional structures. Thus the display records the supplied conjectural schema; it does not supply a uniquely determined categorical conjecture. No quantifier over specializations of \(q\) can be recovered from the record. The statement assumes a semisimple complex group, whereas the supplied literature assessment refers to a general reductive-group formulation; these scopes are retained as distinct.
\subsection{Short English statement}
The source conjectures that a K-theoretic convolution category on an affine Grassmannian is equivalent to a full subcategory of quantum-group-equivariant modules over a quantum coordinate algebra. The group underlying the Grassmannian, the quantum parameter, the categorical constructions, and the intended subcategory require clarification. Supplied dated curation labels the problem partially solved and reports an \(\mathrm{SL}_n\) case as proved; this draft does not independently verify that assessment.
\subsection{Sources}
\begin{itemize}
\item Local UnsolvedMath registry record: permanent ID \(30002885\), title ``K-Theoretic Geometric Satake Equivalence,'' source problem code \texttt{catalogue entry}, in the Oberwolfach Reports Open Problems collection. The supplied record gives proposed year 2015 and names no proposer. Registry provenance: \texttt{data/unsolvedmath/problems.json}, with \texttt{catalog/problems.json} as the supplementary project registry; file membership has not been independently checked.
\item Supplied report citation: \emph{Enveloping Algebras and Geometric Representation Theory} (2016), Oberwolfach Reports, DOI 10.4171/owr/2015/25, pp. 1385--1447. \url{https://doi.org/10.4171/owr/2015/25}. The record describes Oberwolfach Reports as documenting mathematical workshops and their open-problem discussions. The citation's year 2016 is retained alongside the proposed year 2015 and the DOI containing 2015.
\item Supplied literature reference, with attribution as recorded: Kamnitzer--Webster--Weekes--Yacoby, \emph{Quantum K-theoretic geometric Satake}. \url{https://arxiv.org/abs/1509.00112}. The attribution and correspondence between label and link have not been independently verified.
\item Supplied literature reference: \emph{Quantum K-theoretic geometric Satake, \(\mathrm{SL}_n\) case}. \url{https://doi.org/10.1112/S0010437X17007564}. The correspondence between this label and the linked publication has not been independently verified.
\item Literature assessment dated exactly \texttt{2026-08-21}: supplied dated curation, not independently verified. It labels the problem ``partially solved'' and states that the K-theoretic geometric Satake equivalence conjecture is proved for \(\mathrm{SL}_n\), while the stated general reductive-group equivalence remains open. The record characterizes this as literature triage derived from supplied OWR curation and explicitly requests independent verification of time-sensitive status and cited claims.
\end{itemize}
\end{document}
